\documentclass[preprint,12pt,a4paper]{elsarticle}

% ---------- Packages ----------
\usepackage[utf8]{inputenc}
\usepackage[T1]{fontenc}
\usepackage{amsmath,amssymb,amsfonts}
\usepackage{graphicx}
\usepackage{booktabs}
\usepackage{multirow}
\usepackage{array}
\usepackage{xcolor}
\usepackage{hyperref}
\usepackage{cleveref}
\usepackage{algorithm}
\usepackage{algpseudocode}
\usepackage{subcaption}
\usepackage{natbib}

% ---------- Journal info ----------
% \journal{Pattern Recognition}

\begin{document}

% ---------- Title & authors ----------
\begin{frontmatter}

\title{A Multi-Task CNN-Transformer Fusion Network for Authenticity Assessment, Pattern Classification and Defect Detection of Blue Calico}

% Replace with your information
\author[1]{Your Name\corref{cor1}}
\ead{your.email@example.com}
\author[1]{Co-author Name}
\author[2]{Another Co-author}

\affiliation[1]{organization={Your University, Department of Computer Science},
                city={City},
                country={Country}}
\affiliation[2]{organization={Collaborative Institution},
                city={City},
                country={Country}}

\cortext[cor1]{Corresponding author}

\begin{abstract}
Traditional handmade textiles, such as blue calico and embroidery, are valuable cultural heritages whose authenticity assessment, pattern classification and defect detection have long relied on expert experience, making the process subjective, time-consuming and difficult to scale.
Existing automated methods mostly tackle these tasks independently and rely solely on convolutional neural networks, which are good at local texture but weak at capturing global pattern structures.
To address these limitations, we propose a unified multi-task framework, named \textbf{XXXNet}, that fuses a CNN branch for fine-grained texture with a DINOv2-based Vision Transformer branch for global semantic structure.
A lightweight Patch Transformer is further introduced to aggregate spatial relationships among patch tokens.
For authenticity assessment, we employ an ArcFace head to enlarge the decision boundary between genuine and forged samples.
A dynamic loss weighting strategy is adopted to balance the three heterogeneous tasks during training.
Extensive experiments on the [Dataset Name] dataset demonstrate that our method achieves [xx.x\% / xx.x\% / xx.x\%] accuracy on authenticity, pattern and defect classification, respectively, outperforming several representative CNN and Transformer baselines.
Ablation studies validate the effectiveness of each proposed component.
\end{abstract}

\begin{highlights}
\item A unified multi-task network is proposed for authenticity, pattern and defect analysis of traditional textiles.
\item CNN and self-supervised Vision Transformer features are fused to model both local texture and global structure.
\item ArcFace loss and dynamic task weighting are introduced to improve authenticity discrimination and training stability.
\item Comprehensive experiments and ablation studies demonstrate the superiority of the proposed design.
\end{highlights}

\begin{keywords}
Cultural heritage digitization \sep blue calico authentication \sep multi-task learning \sep Vision Transformer \sep self-supervised learning \sep ArcFace
\end{keywords}

\end{frontmatter}

% ---------- 1. Introduction ----------
\section{Introduction}\label{sec:intro}

Blue calico (Lan Yin Hua Bu) is one of the representative traditional handicrafts in China.
Its authentication, pattern categorization and quality inspection are essential for cultural heritage protection and market regulation.
Currently, these tasks are mainly performed by experienced craftsmen or experts, which is not only labor-intensive but also subjective and inconsistent across different evaluators \citep{sample-ref}.
Therefore, an automated, objective and efficient intelligent assessment system is highly desirable.

In recent years, deep learning has shown promising results in various computer vision tasks.
Convolutional neural networks (CNNs), such as ResNet \citep{he2016deep} and EfficientNet \citep{tan2019efficientnet}, have been widely adopted for texture classification and defect detection.
However, CNNs are inherently limited in modeling long-range dependencies and global pattern layouts, which are crucial for distinguishing subtle texture forgeries and understanding traditional textile patterns.

Meanwhile, Vision Transformers (ViTs) and self-supervised foundation models such as DINOv2 \citep{oquab2023dinov2} have demonstrated strong capability in learning semantic visual representations.
Nevertheless, applying pure Transformer backbones directly may overlook the fine-grained local cues that are important for defect localization and subtle authenticity cues.

Motivated by the above observations, this paper proposes a multi-task CNN-Transformer fusion network that integrates the complementary strengths of both architectures.
Specifically, a ResNet50 branch extracts texture and local details, while a DINOv2-S/14 branch with LoRA fine-tuning captures global pattern semantics.
The features from both branches are fused and fed into three task-specific heads.
For authenticity assessment, we replace the conventional linear head with an ArcFace head \citep{deng2019arcface} to enforce a large angular margin between genuine and forged samples.
Moreover, a dynamic loss weighting strategy is employed to balance the three tasks with different loss scales.

The main contributions of this work are summarized as follows:
\begin{itemize}
    \item We propose a unified multi-task framework for simultaneous authenticity assessment, pattern classification and defect detection of blue calico.
    \item We design a CNN-Transformer fusion architecture that leverages ResNet50 for local texture and DINOv2 for global pattern structure, enhanced by a Patch Transformer for spatial aggregation.
    \item We introduce an ArcFace-based authenticity head and a dynamic loss weighting strategy to improve discriminability and training stability.
    \item Extensive experiments demonstrate that the proposed method outperforms several strong CNN and Transformer baselines, and ablation studies verify the effectiveness of each component.
\end{itemize}

% ---------- 2. Related Work ----------
\section{Related Work}\label{sec:related}

\subsection{Traditional Textile Authentication and Quality Inspection}

Early works on textile authentication mainly relied on hand-crafted features, such as color histograms, local binary patterns (LBP), and GLCM, followed by classical classifiers like SVM or k-NN.
With the development of deep learning, CNN-based methods have become dominant.
For example, \citet{sample-ref} employed MobileNet for embroidery authenticity classification.
However, most existing works treat authentication, pattern classification and defect detection as independent tasks, ignoring their inherent correlations.

\subsection{Multi-Task Learning in Computer Vision}

Multi-task learning (MTL) aims to improve generalization by sharing representations across related tasks.
It has been successfully applied in face analysis, scene understanding and medical imaging.
In the context of cultural heritage, MTL can naturally model the correlations among authenticity, pattern and defect tasks: genuine textiles usually exhibit regular patterns without defects, while forgeries may contain both pattern inconsistencies and local anomalies.

\subsection{Vision Transformers and Self-Supervised Learning}

Vision Transformers \citep{dosovitskiy2020image} have achieved competitive or superior performance compared with CNNs on many image recognition benchmarks.
Subsequent works such as Swin Transformer \citep{liu2021swin} introduced hierarchical structures and shifted windows to improve efficiency.
More recently, self-supervised foundation models like DINOv2 \citep{oquab2023dinov2} have shown strong transferability by pre-training on large-scale unlabeled images.
Parameter-efficient fine-tuning methods such as LoRA \citep{hu2021lora} make it feasible to adapt these large models to domain-specific datasets.

\subsection{Deep Metric Learning}

ArcFace \citep{deng2019arcface} is a popular additive angular margin loss for face recognition.
It has also been applied to fine-grained classification, anomaly detection and open-set recognition.
Its core idea is to enforce a margin between different classes in the angular space, leading to more discriminative embeddings.
In blue calico authentication, where genuine and forged samples can be visually subtle, ArcFace is well suited to separate the two classes.

% ---------- 3. Method ----------
\section{Method}\label{sec:method}

\subsection{Problem Formulation}

Given an input image $x$, the model is trained to predict three labels simultaneously:
\begin{itemize}
    \item authenticity label $y_a \in \{0, 1\}$ (0: forged, 1: genuine);
    \item pattern label $y_p \in \{1, \dots, K\}$;
    \item defect label $y_d \in \{0, 1\}$ (0: no defect, 1: defect).
\end{itemize}
The overall objective is to minimize a multi-task loss:
\begin{equation}
    \mathcal{L}_{total} = w_a \mathcal{L}_a + w_p \mathcal{L}_p + w_d \mathcal{L}_d,
\end{equation}
where $w_a$, $w_p$ and $w_d$ are task weights that are dynamically adjusted during training.

\subsection{Network Architecture}

The proposed network consists of three parts: a CNN branch, a Vision Transformer branch and a fusion module with task-specific heads.
An overview is shown in \cref{fig:architecture}.

\begin{figure}[htbp]
    \centering
    \includegraphics[width=0.95\linewidth]{figures/network_architecture.pdf}
    \caption{Overview of the proposed multi-task CNN-Transformer fusion network. The input image is processed by a ResNet50 branch and a DINOv2-S/14 branch. Patch tokens from DINOv2 are further aggregated by a Patch Transformer. The fused features are fed into three task-specific heads, where the authenticity head uses ArcFace loss.}
    \label{fig:architecture}
\end{figure}

\subsubsection{CNN Branch}

The CNN branch is based on ResNet50 pretrained on ImageNet.
We remove the final classification layer and use the output of the global average pooling layer as the texture feature vector $\mathbf{f}_{cnn} \in \mathbb{R}^{2048}$.

\subsubsection{Vision Transformer Branch}

The Transformer branch uses DINOv2-S/14 as the backbone.
We add LoRA adapters to the query-key-value projection layers for parameter-efficient fine-tuning.
The class token is taken as the global semantic feature $\mathbf{f}_{cls} \in \mathbb{R}^{384}$.
In addition, the patch tokens are passed through a lightweight Transformer encoder (Patch Transformer) to obtain a spatially aggregated feature $\mathbf{f}_{patch} \in \mathbb{R}^{384}$.

\subsubsection{Feature Fusion and Heads}

The three features are concatenated and projected to a compact representation:
\begin{equation}
    \mathbf{f} = \text{MLP}([\mathbf{f}_{cnn}; \mathbf{f}_{cls}; \mathbf{f}_{patch}]).
\end{equation}

For pattern classification and defect detection, we use standard linear heads with cross-entropy loss:
\begin{align}
    \mathcal{L}_p &= \text{CrossEntropy}(\hat{y}_p, y_p), \\
    \mathcal{L}_d &= \text{CrossEntropy}(\hat{y}_d, y_d).
\end{align}

For authenticity assessment, we use an ArcFace head.
Let $\mathbf{e}$ be the feature embedding and $W$ be the weight matrix of the last fully connected layer.
ArcFace adds an angular margin $m$ to the target class:
\begin{equation}
    \mathcal{L}_a = -\log \frac{e^{s \cos(\theta_{y_a} + m)}}{e^{s \cos(\theta_{y_a} + m)} + \sum_{j \neq y_a} e^{s \cos \theta_j}},
\end{equation}
where $\theta_j$ is the angle between $\mathbf{e}$ and the $j$-th column of $W$, and $s$ is a scaling factor.

\subsection{Dynamic Loss Weighting}

Because the three tasks have different loss magnitudes and convergence speeds, we adopt a dynamic weighting strategy:
\begin{equation}
    w_i(t) = \frac{\mathcal{L}_i(t-1)}{\mathcal{L}_a(t-1) + \mathcal{L}_p(t-1) + \mathcal{L}_d(t-1)}, \quad i \in \{a, p, d\}.
\end{equation}
Intuitively, tasks with higher loss receive larger weights, preventing any single task from dominating the gradient updates.

% ---------- 4. Experiments ----------
\section{Experiments}\label{sec:experiments}

\subsection{Dataset and Evaluation Metrics}

\textbf{Dataset.}
We evaluate the proposed method on the [Dataset Name] dataset, which contains [number] images of blue calico textiles.
The dataset is divided into [training/validation/test] splits with [xxx / xxx / xxx] images.
Each image is annotated with three labels: authenticity, pattern category and defect presence.

\textbf{Metrics.}
We report accuracy and AUC for authenticity and defect classification, and accuracy for pattern classification.
All experiments are repeated [X] times and the mean results are reported.

\subsection{Implementation Details}

All images are resized to $224 \times 224$ pixels.
Data augmentation includes random horizontal flipping, rotation and color jittering.
We use AdamW as the optimizer with a cosine learning rate schedule.
The backbone learning rate is set to $1 \times 10^{-4}$ and the head learning rate to $1 \times 10^{-3}$.
The batch size is 32 and the model is trained for 50 epochs unless otherwise stated.
All experiments are conducted on an NVIDIA RTX 4060 Ti 16 GB GPU.

\subsection{Comparison with State-of-the-Art Methods}

\Cref{tab:main-result} compares the proposed method with representative CNN and Transformer baselines.

\begin{table}[htbp]
    \centering
    \caption{Comparison with external baselines on the corrected test set. The best and second-best results are highlighted in \textbf{bold} and \underline{underlined}, respectively.}
    \label{tab:main-result}
    \begin{tabular}{lccc}
        \toprule
        Method & Auth Acc ($\%$) & Pat Acc ($\%$) & Def Acc ($\%$) \\
        \midrule
        ResNet50 & 95.9 & 98.4 & \textbf{94.0} \\
        EfficientNet-B3 & 93.8 & 99.6 & 92.8 \\
        EfficientNet-B4 & 94.6 & 100.0 & 92.1 \\
        ConvNeXt-Tiny & 90.5 & 99.2 & 92.3 \\
        DINOv2-S/14 + LoRA & 88.0 & 98.4 & 92.8 \\
        \midrule
        Ours (XXXNet) & \textbf{xx.x} & \underline{xx.x} & \underline{xx.x} \\
        \bottomrule
    \end{tabular}
\end{table}

From \cref{tab:main-result}, we observe that [add your analysis here].

\subsection{Ablation Studies}

To validate the contribution of each component, we conduct ablation experiments as shown in \cref{tab:ablation}.

\begin{table}[htbp]
    \centering
    \caption{Ablation study on the proposed components.}
    \label{tab:ablation}
    \begin{tabular}{lccc}
        \toprule
        Variant & Auth Acc ($\%$) & Pat Acc ($\%$) & Def Acc ($\%$) \\
        \midrule
        Ours full & xx.x & xx.x & xx.x \\
        w/o Patch Transformer & 96.1 & 99.7 & 92.5 \\
        w/o ArcFace & 97.0 & 98.4 & 92.0 \\
        w/o DINOv2 branch & xx.x & xx.x & xx.x \\
        w/o ResNet50 branch & xx.x & xx.x & xx.x \\
        \bottomrule
    \end{tabular}
\end{table}

The results indicate that [explain].

\subsection{Visualization and Case Study}

We visualize the regions that the model focuses on using Grad-CAM.
As shown in \cref{fig:cam}, the authenticity head concentrates on subtle texture and edge inconsistencies, while the defect head focuses on local abnormal regions.
These visualizations are consistent with human expert judgment, demonstrating the interpretability of the proposed method.

\begin{figure}[htbp]
    \centering
    \includegraphics[width=0.95\linewidth]{figures/gradcam_examples.pdf}
    \caption{Grad-CAM visualization of the authenticity head and defect head. Warmer colors indicate higher attention weights.}
    \label{fig:cam}
\end{figure}

% ---------- 5. Conclusion ----------
\section{Conclusion}\label{sec:conclusion}

In this paper, we propose a multi-task CNN-Transformer fusion network for blue calico authentication, pattern classification and defect detection.
The proposed method effectively combines the local texture representation of CNNs and the global semantic understanding of self-supervised Vision Transformers.
ArcFace loss and dynamic loss weighting further improve the authenticity discrimination and training stability.
Experimental results demonstrate the superiority of the proposed method over several strong baselines.

For future work, we plan to extend the defect classification to pixel-level defect segmentation, explore lightweight deployment on mobile devices, and validate the method on more diverse cross-dataset scenarios.

% ---------- Acknowledgments ----------
\section*{Acknowledgments}

This work was supported by [Funding Source, Grant Number].
We thank the reviewers for their constructive comments.

% ---------- References ----------
\bibliographystyle{elsarticle-num}
\bibliography{references}

\end{document}
